\documentclass[11pt]{article}
\usepackage[margin=1in]{geometry}
\usepackage[T1]{fontenc}
\usepackage{lmodern}
\usepackage{microtype}
\usepackage{amsmath,amssymb}
\usepackage{booktabs}
\usepackage{float}
\usepackage{multirow}
\usepackage{enumitem}
\usepackage{xcolor}
\usepackage[colorlinks=true,linkcolor=blue!50!black,citecolor=blue!50!black,urlcolor=blue!50!black]{hyperref}
\setlist{itemsep=2pt}

% Numbers still to be filled in from results/<dataset>-<split>/summary.md and logs/check_data.log
\newcommand{\tbd}{\textcolor{red}{\textbf{??}}}
\newcommand{\bok}{best-of-$K$}

\title{Five Multi-Solution Benchmarks and a Prior-Only Baseline\\[4pt]
\large Project report v2: datasets, baseline, reference numbers}
\author{Harshit Bisht (2025AIZ8455), Ashwendh Choudhary (2026AIY7513)}
\date{November 2026}

\begin{document}
\maketitle

\begin{abstract}
We study GRAM~\cite{baek2026gram}, a recursive reasoning model with a latent variable, on problems with many correct
answers, around three questions: how to choose its KL weight $\beta$ (Q1), whether training the latent through an
answer-seeing posterior is ever harmful (Q2), and whether a model produces valid solutions it never saw in training
(Q3). Before testing any intervention we fix the ground: five datasets whose solution sets are enumerated exactly
(mazes, N-Queens, Sudoku, graph 3-colouring and random 3-SAT), each storing at most 16 solutions per training input so
that the rest are held out; one baseline that never feeds the answer to the latent (GRAM's network trained with
$\mathcal{N}(0,I)$ latents and a best-of-$K$ objective); and one set of metrics computed by the same code on every
task. This report describes the three and records the baseline's numbers. Everything later is measured against
them.
\end{abstract}

\section{Why this report}

Earlier work in the project (report v1) compared deterministic TRMs, one-of-many objectives and GRAM on mazes and
N-Queens, with metrics defined slightly differently on each task. Following the supervisor's redirection (6 Nov), the
project now studies GRAM itself (Q1--Q3). Answers to these questions are comparisons between models and between
tasks, so they need datasets built the same way, a baseline trained the same way on all of them, and metrics that
mean the same thing everywhere. That is the content of this report; it contains no intervention.

\section{Datasets}

\subsection{Common construction}

All five datasets are built by \texttt{scripts/build\_multi\_datasets.sh} and share one design.
\begin{itemize}
  \item \textbf{Exact solution sets.} Every solution of every input is enumerated (dynamic programming over BFS
  layers for mazes, backtracking or bitmask search elsewhere), so the solution count $k(x)$ is exact.
  \item \textbf{Controlled ambiguity.} Each task has a knob that moves $k(x)$ from 1 to hundreds or thousands
  (Table~\ref{tab:datasets}). Inputs are drawn until each group of solution counts has its quota: one third have a
  unique solution, the rest are spread evenly over groups (for Sudoku, graph colouring and SAT: 2--3, 4--15, 16--63,
  64--255, 256+; mazes use log$_2$-spaced groups). So every level of ambiguity is equally represented, rather than
  whatever the generator produces most often. N-Queens is the exception: its inputs are all boards with 1--3 queens
  given, with their natural distribution of counts.
  \item \textbf{Stored and held-out solutions.} Each training input stores at most 16 of its solutions, drawn
  uniformly; training only ever sees these. On inputs with $k(x)>16$ the other $k(x)-16$ are held out, which is what
  Q3 measures. The exact count is kept with every input.
  \item \textbf{Disjoint splits.} Test inputs never appear in training (checked; Section~\ref{sec:checks}). This is a
  change from report v1, where N-Queens test boards were also training boards.
  \item \textbf{Variants.} Four training sets share one test split: all stored solutions (naive training), one row per
  input with its set of stored solutions (\texttt{-1oml}, used by the baseline), one fixed solution per input
  (\texttt{-single}) and unique-solution inputs only (\texttt{-unique}).
\end{itemize}

\subsection{The five tasks}

\begin{table}[h]
\centering\small
\begin{tabular}{@{}lllllll@{}}
\toprule
Task & Input & Ambiguity knob & Solutions & Train / test & Seq.\ len & Valid if \\
\midrule
Maze & 15$\times$15 grid & walls removed & 1--924 & 100k / 2k & 225 & a shortest path \\
N-Queens & 10$\times$10 board & queens given & 1--93 & 36{,}907 / 6{,}513 & 100 & no attacks, givens kept \\
Sudoku & 9$\times$9 grid & clues & 1--1024 & 100k / 2k & 81 & valid, clues kept \\
Graph col. & 10 nodes & edge density & 1--9842 & 100k / 2k & 55 & proper 3-colouring \\
3-SAT & 16 variables & clause density & 1--10{,}576 & 100k / 2k & 232 & all clauses true \\
\bottomrule
\end{tabular}
\caption{The five datasets. Solution ranges are over train and test (graph colouring: up to 6561 in train, 9842 in
test).}
\label{tab:datasets}
\end{table}

\paragraph{Mazes.} A spanning-tree maze on a 7$\times$7 lattice of rooms, with each remaining inner wall removed with
some probability to create loops, and a start and goal at least 8 steps apart. The label marks the cells of one
shortest path; any shortest path is correct. The number of shortest paths and each cell's share of them are computed
exactly.

\paragraph{N-Queens.} All 724 solutions of 10-Queens, with 7, 8 or 9 queens removed; each distinct partial board is
one input, and its completions are its solutions. 15\% of distinct inputs form the test split. Training inputs store
up to 16 completions; test inputs store all of them.

\paragraph{Sudoku.} A random complete grid with 26--50 clues kept at random cells; puzzles with more than 1024
solutions are discarded. The label is the whole solved grid.

\paragraph{Graph 3-colouring.} A colouring of 10 nodes is planted and each pair of differently coloured nodes is joined
with probability $p\sim U[0.1, 1]$, so the graph is always 3-colourable. Colourings are counted in \emph{canonical
form} (colours numbered by first appearance), so relabelling the colours never makes a new solution: two solutions
differ in which nodes share a colour. The maximum, 9842, is the number of partitions of 10 nodes into at most 3
groups (the empty graph). The input is the upper triangle of the adjacency matrix; the answer is the first 10 tokens.
Because only about a thousand graphs have $\geq$4096 colourings, the top group is 1024--9842.

\paragraph{Random 3-SAT.} 16 variables and $\alpha\cdot 16$ random 3-clauses with $\alpha\sim U[1, 4.5]$; up to a
quarter of the variables may be given (fixed to their value in a random solution). Unsatisfiable formulas are
dropped. Few clauses leave thousands of satisfying assignments; near the threshold ($\alpha\approx4.27$) a handful.
The answer is the 16 variable slots.

\subsection{Checks}
\label{sec:checks}

\texttt{scripts/analyze\_multi\_samples.py --check\_data} checks each dataset with the same checkers the evaluation
uses: every stored solution valid and distinct, the number stored equal to $\min(k(x), 16)$ on training inputs, and
no input shared between train and test. On small builds of all five tasks every check passes, and a deliberately
corrupted copy (one broken solution, one duplicate, one leaked input) fails all three corresponding checks. On the full
datasets every check passes on every input of both splits for mazes, N-Queens, Sudoku and graph colouring; for SAT on all test inputs and the first 20k training inputs. Table~\ref{tab:groups} gives the inputs per group of solution counts.

\begin{table}[h]
\centering\small
\begin{tabular}{@{}lrrrrrrr@{}}
\toprule
& \multicolumn{6}{c}{Solutions per input} & Stored \\
\cmidrule(lr){2-7}
Dataset, split & 1 & 2--3 & 4--15 & 16--63 & 64--255 & 256+ & solutions \\
\midrule
Maze, train & 33{,}333 & 7{,}408 & 14{,}816 & 14{,}815 & 14{,}814 & 14{,}814 & 885{,}729 \\
Maze, test & 667 & 149 & 296 & 296 & 296 & 296 & 17{,}701 \\
N-Queens, train & 13{,}586 & 15{,}427 & 7{,}534 & 293 & 67 & 0 & 102{,}513 \\
N-Queens, test & 2{,}398 & 2{,}705 & 1{,}358 & 43 & 9 & 0 & 18{,}466 (all) \\
Sudoku, train & 33{,}333 & 13{,}334 & 13{,}334 & 13{,}333 & 13{,}333 & 13{,}333 & 800{,}363 \\
Sudoku, test & 667 & 267 & 267 & 267 & 266 & 266 & 15{,}920 \\
Graph col., train & 33{,}333 & 11{,}112 & 11{,}111 & 11{,}111 & 11{,}111 & 22{,}222 & 850{,}425 \\
Graph col., test & 667 & 223 & 222 & 222 & 222 & 444 & 17{,}021 \\
3-SAT, train (first 20k) & 6{,}584 & 2{,}597 & 2{,}659 & 2{,}709 & 2{,}700 & 2{,}751 & 165{,}067 \\
3-SAT, test & 667 & 267 & 267 & 267 & 266 & 266 & 16{,}385 \\
\bottomrule
\end{tabular}
\caption{Inputs per group of solution counts. Mazes are drawn to log$_2$-spaced groups (2--3, 4--7, 8--15, \dots),
so 2--3 holds one group and 4--15 two; graph colouring's top two groups (256--1023, 1024--9842) both fall in 256+.
N-Queens test inputs store all their solutions.}
\label{tab:groups}
\end{table}

\paragraph{Consequences.} N-Queens has its natural distribution of counts: only 360 training inputs (1\%) have more
than 16 solutions, none more than 93, so Q3 on N-Queens rests on few inputs and small held-out sets. Mazes, Sudoku,
graph colouring and SAT each have about 40--44k training inputs with 16 or more solutions. N-Queens also has a third as many training inputs as
the others, so at the same step budget it is seen about 130 times instead of 48.

\section{Baseline: prior-only best-of-$K$}

\subsection{Model}

The network is GRAM's~\cite{baek2026gram}, itself a TRM~\cite{jolicoeur2025trm} with a latent variable: a 2-layer
transformer block (hidden size 512, 8 heads, $\sim$10.9M parameters) applied recursively, with 3 high-level cycles of
4 low-level cycles per supervision step and up to 16 supervision steps (deep supervision with a learned halting head).
After each high-level cycle, every output position $i$ gets a fresh latent $z_i\in\mathbb{R}^{256}$, projected and
added to the state: $h_i \leftarrow h_i + W z_i + b$. A rollout therefore draws a new latent up to $3\times16=48$
times.

GRAM draws $z$ during training from a posterior $q(z\mid x, y)$ that sees the target, at test time from a learned prior
$p(z\mid x)$, and adds $\beta\,\mathrm{KL}(q\,\|\,p)$ to the loss. The baseline removes the posterior: the latent never
sees the answer, during training or test, so there is no KL term and no $\beta$. This is the simplest model that can
answer Q2 (can we do without the posterior?) and gives Q1 a reference that has no $\beta$ at all.

\paragraph{Noise distribution.} The baseline draws $z_i\sim\mathcal{N}(0,I)$. Because $z$ enters through the learned
map $W$, the injected noise $Wz_i+b$ is $\mathcal{N}(b, WW^\top)$: its mean and covariance are still learned, but they
are the same for every input and position, and nothing in the loss can make them diverge. We first used GRAM's learned
prior $p(z\mid x)$ instead, whose log-variance is unconstrained once the KL term is gone. That version went NaN late in
training on graph colouring, and on N-Queens its training accuracy fell from about 0.7 to 0.2 over the last fifth of
training, while the $\mathcal{N}(0,I)$ version trained smoothly on both. We keep the learned-prior runs as an ablation
(Section~\ref{sec:ablation}).

\subsection{Objective}

Each training input $x$ with stored solutions $S(x)$ is rolled out $K=8$ times with independent latents. At every
supervision step, rollout $j$ produces token distributions $f_j$, and
\begin{equation}
  \mathcal{L}(x) = \min_{j\le K}\ \min_{s\in S(x)}\ \ell(f_j, s)
  \;+\; \frac1K\sum_{j\le K} \mathcal{L}_{\text{halt}}(f_j),
\end{equation}
where $\ell$ is the mean token cross-entropy (stablemax, as in TRM) and the halting loss trains each rollout's halting
head to predict whether its current answer is valid. Only the best rollout receives the reconstruction gradient,
towards its closest stored solution; it is weighted by $K$ so that the loss has the same scale as with one rollout.
The best rollout is chosen afresh at every supervision step. This is best-of-many sample
training~\cite{bhattacharyya2018bestofmany}, here combined with MinLoss over solutions~\cite{nandwani2021oneofmany}.

\paragraph{What this objective rewards.} Any model whose every sample is some stored solution reaches zero loss,
whether it always returns the same one or spreads evenly over all of them. The objective is therefore indifferent
between collapse and coverage; whatever diversity the baseline shows comes from training dynamics, not from the loss.
This is a deliberate choice for a baseline: it is the weakest diversity pressure that still trains, so any
intervention that targets diversity has a clean reference. Variants that do reward coverage (IMLE~\cite{li2018imle},
which picks a random stored target and trains the nearest sample to it; matching samples to solutions one to one) are
interventions, not baselines.

\paragraph{Halting target.} The halting head's target is validity under the task's checker for mazes, N-Queens,
Sudoku and graph colouring; for SAT it is ``matches the closest stored solution'', which undercounts valid answers.
\texttt{train/exact\_accuracy} in the training logs follows the same definitions and is a single-sample rate over all
rollouts, not best-of-$K$.

\subsection{Training}

Batch 768 sequences ($96$ inputs $\times$ 8 rollouts), learning rate $10^{-4}$ (constant after warm-up), weight decay
1.0, EMA weights, $\sim$50k steps, so each 100k-input dataset is seen about 48 times. Rollouts of one input halt
together. SAT uses learned position embeddings instead of RoPE, so each clause slot has a fixed embedding. One seed
(42). Launch: \texttt{bash scripts/multi\_compare.sh <dataset> gram\_std\_bestofk}.

\section{Evaluation protocol and metrics}

Each model draws 100 samples per input, each with fresh latents, running all 16 supervision steps
(\texttt{scripts/multi\_evaluate.sh}). The test split has 2000 inputs; for Q3 we also sample the first 2000 training
inputs (training inputs are stored in random order). All metrics come from \texttt{scripts/analyze\_multi\_samples.py}
and are reported for all inputs and per group of solution counts (1, 2--3, 4--15, 16--63, 64--255, 256+).

\begin{itemize}
  \item \textbf{Validity}: the share of samples the task's checker accepts, i.e.\ expected single-sample accuracy.
  \textbf{Any-valid@$K$}: an input counts if any of its first $K$ samples is valid (an upper bound on what a verifier
  could pick).
  \item \textbf{Distinct valid@$K$} on multi-solution inputs, and the \textbf{ratio to uniform}: distinct valid
  answers divided by the number a uniform sampler over the input's $k$ solutions would find with the same number $v$
  of valid draws, $k\,(1-(1-1/k)^{v})$. 1 means as diverse as uniform sampling, $1/v$ means one solution repeated. The
  ratio separates diversity from validity. \textbf{Coverage}: distinct valid / $k$.
  \item \textbf{Unique-solution inputs}: distinct answers per input (ideally 1: noise should not change an answer that
  is forced).
  \item \textbf{Error anatomy} of invalid samples: cells differing from the nearest stored solution, and on
  multi-solution inputs the share whose every cell agrees with some stored solution while matching none (a
  \emph{blend} of solutions rather than an error of its own). Both are bounds, since only 16 solutions are stored.
  \item \textbf{Held-out solutions} (training split, inputs with $k>16$): the share of valid samples that are not among
  the stored ones, against a uniform sampler's $1-16/k$. Near 0: the model reproduces what it was trained on. Near
  the uniform share: it produces unseen solutions as readily as seen ones. We also report the share of unseen valid
  samples that recombine stored ones cell by cell; this must be read against a uniform sampler's value, which is high
  (about 97\% for graph colouring).
  \item \textbf{Graph colouring only}: validity up to relabelling, which accepts proper colourings whose colours are
  not in canonical order.
\end{itemize}

\paragraph{The metrics on known samplers.} On synthetic samples the analyzer returns the expected values on all five
tasks: a uniform sampler over all solutions has ratio 1.00 and an unseen share equal to the uniform share (83.9\% vs
83.6\%); a collapsed sampler finds 1 solution and 0\% unseen; a sampler uniform over the stored solutions has ratio 1
while $k\le16$ and less above, and 0\% unseen; blends of stored solutions are recognised as blends (98--100\%) and
random noise is not (0--2\%); relabelled colourings are 0\% valid strictly and 100\% up to relabelling.

\section{Baseline results at 50k steps}

\paragraph{Status.} These are the first, 50k-step versions of the baseline, recorded before any longer run. They are
a snapshot, not the final reference: we have not checked convergence, and the maze model is the likeliest to be
undertrained (in report v1, maze models kept improving between 50k and 100k steps). The reference baseline will be
trained for 150k steps with two seeds, with every checkpoint scored on a validation split (Section~\ref{sec:next}).
3-SAT is reported but did not learn (Section~\ref{sec:obs}).

\begin{table}[h]
\centering\small
\begin{tabular}{@{}lccccc@{}}
\toprule
& Maze & N-Queens & Sudoku & Graph col. & 3-SAT \\
\midrule
Validity, all inputs & 76.7 & 96.9 & 99.7 & 99.5 & 0.4 \\
\quad unique-solution inputs & 75.0 & 93.2 & 99.8 & 100.0 & 0.0 \\
\quad multi-solution inputs & 77.6 & 99.1 & 99.6 & 99.3 & 0.7 \\
Any-valid@100 & 93.2 & 100.0 & 100.0 & 100.0 & 12.9 \\
Distinct valid@100 (multi) & 1.7 & 3.3 & 4.6 & 35.7 & 0.6 \\
Ratio to uniform & 0.13 & 0.89 & 0.27 & 0.75 & n/a \\
Coverage & 9.9 & 89.2 & 24.0 & 46.8 & 0.5 \\
Distinct answers, unique inputs & 13.6 & 5.17 & 1.11 & 1.02 & 93.7 \\
Invalid: cells off nearest stored (median) & 13 & 7 & 22 & 1 & 6 \\
Invalid: blends (multi) & 60.7 & 22.7 & 3.3 & 94.5 & 22.0 \\
\midrule
\multicolumn{6}{@{}l}{\emph{Training inputs with more than 16 solutions (Q3)}} \\
Inputs & 881 & 14 & 801 & 876 & 781 \\
Unseen share of valid samples & 72.0 & 30.6 & 76.9 & 84.7 & n/a \\
\quad uniform sampler & 71.9 & 35.7 & 76.8 & 84.7 & n/a \\
\bottomrule
\end{tabular}
\caption{Prior-only best-of-8 with $\mathcal{N}(0,I)$ latents (\texttt{gram\_std\_bestofk}, seed 42, 50k steps), 100
samples per input, test split (N-Queens: 6513 inputs, others 2000) except the Q3 rows (first 2000 training inputs).
Multi-solution validity is the input-weighted mean over groups. n/a: 3-SAT has too few valid samples (0.4\%) for
diversity or held-out shares to mean anything. Diversity averages differ across tasks partly because
the tasks have different mixes of solution counts; Table~\ref{tab:ratio} compares them group by group. Per-group
tables are in Appendix~\ref{app:groups}.}
\label{tab:baseline}
\end{table}

\begin{table}[h]
\centering\small
\begin{tabular}{@{}lccccc@{}}
\toprule
Solutions & 2--3 & 4--15 & 16--63 & 64--255 & 256+ \\
\midrule
Maze & 0.48 & 0.17 & 0.06 & 0.05 & 0.10 \\
N-Queens & 0.92 & 0.84 & 0.71 & 0.61 & -- \\
Sudoku & 0.61 & 0.37 & 0.15 & 0.09 & 0.10 \\
Graph colouring & 0.94 & 0.81 & 0.59 & 0.58 & 0.80 \\
3-SAT & n/a & n/a & n/a & n/a & n/a \\
\bottomrule
\end{tabular}
\caption{Ratio of distinct valid answers to a uniform sampler's, by number of solutions ($\mathcal{N}(0,I)$, 50k
steps, test split). N-Queens has 43 test inputs at 16--63 and 9 at 64--255, and none above 93 solutions.}
\label{tab:ratio}
\end{table}

\subsection{First observations}
\label{sec:obs}

These describe the numbers; they are not yet claims, since every run is one seed at 50k steps.
\begin{itemize}
  \item \textbf{3-SAT did not learn.} 0.4\% of samples are valid, none on unique-solution formulas, where 100 samples
  give 93.7 distinct answers: the outputs are close to random. Its training loss rose over the run, and in October
  GRAM failed on SAT the same way, so the cause is likely in the task setup rather than this run; it needs its own
  diagnosis before SAT can serve as a benchmark. Even its best group (256+ solutions, 2.5\% valid) may be no better
  than a random assignment, which is valid with probability $k/2^{\text{free variables}}\ge 0.4\%$ there and up to
  about 11\% (the test maximum, 7012 solutions over $2^{16}$ assignments); the chance baseline (Section~\ref{sec:next}) will settle it. The learned prior fails the same way.
  \item \textbf{Single samples are valid on three of the other four tasks}: 96.9--99.7\% on N-Queens, Sudoku and
  graph colouring. Mazes are the exception (76.7\%), with validity falling from 89\% at 16--63 paths to 53\% at 256+.
  \item \textbf{The same objective gives very different diversity on different tasks}, also at matched numbers of
  solutions (Table~\ref{tab:ratio}). At 16--63 solutions the ratio to uniform is 0.71 on N-Queens and 0.59 on graph
  colouring but 0.15 on Sudoku and 0.06 on mazes. Graph colouring keeps it at 256+ (74.8 distinct valid colourings in
  100 samples, uniform 93.4); Sudoku and mazes fall to about 0.1 (9.1 and 2.7 distinct).
  \item \textbf{Whether noise changes forced answers is also task-dependent.} On unique-solution inputs, Sudoku and
  graph colouring give one answer (1.11 and 1.02 distinct per input). Mazes give 13.6 distinct answers at 75\%
  validity and N-Queens 5.2 at 93\%: there the latent varies errors rather than choosing among solutions.
  \item \textbf{Training inputs are sampled like new ones.} Validity and diversity on the first 2000 training inputs
  match the test split group by group on every task.
  \item \textbf{No preference for stored solutions (Q3).} On training inputs with more than 16 solutions, the share
  of valid samples that were never stored for that input equals a uniform sampler's to within about 1 point on
  mazes, Sudoku and graph colouring (881, 801 and 876 inputs), with both noise types. The baseline produces unseen
  solutions exactly as often as seen ones. This says nothing about \emph{how many} it produces: on mazes, 1.6 distinct
  unseen paths per input. The recombination share (unseen samples made only of cells from stored solutions; 58\% on
  Sudoku, 76\% on mazes, 98\% on graph colouring) cannot be read without a uniform reference.
  \item \textbf{Errors differ in kind}: one cell from a stored solution on graph colouring (median), 7 on N-Queens,
  13 on mazes, 22 on Sudoku.
\end{itemize}

\subsection{Ablation: learned prior}
\label{sec:ablation}

The same objective with GRAM's learned prior $p(z\mid x)$ (\texttt{gram\_prior\_bestofk}), whose variance nothing
constrains once the KL term is gone. Graph colouring went NaN at step 43{,}071, after its halting loss became erratic
from about step 37.5k; we score its last checkpoint before that, step 37{,}494 (the one exception to the
final-checkpoint rule), and also step 41{,}660, inside the erratic phase.

\begin{table}[h]
\centering\small
\begin{tabular}{@{}lcccccc@{}}
\toprule
& \multicolumn{2}{c}{Validity} & \multicolumn{2}{c}{Distinct answers, unique} & \multicolumn{2}{c}{Ratio to uniform} \\
\cmidrule(lr){2-3}\cmidrule(lr){4-5}\cmidrule(l){6-7}
Task & learned & $\mathcal{N}(0,I)$ & learned & $\mathcal{N}(0,I)$ & learned & $\mathcal{N}(0,I)$ \\
\midrule
Maze & 74.2 & 76.7 & 11.7 & 13.6 & 0.15 & 0.13 \\
N-Queens & 48.0 & 96.9 & 37.2 & 5.17 & 0.91 & 0.89 \\
Sudoku & 99.5 & 99.7 & 1.10 & 1.11 & 0.40 & 0.27 \\
Graph col., step 37{,}494 & 37.0 & 99.5 & 7.45 & 1.02 & 0.74 & 0.75 \\
Graph col., step 41{,}660 & 7.1 & & 69.2 & & 0.84 & \\
3-SAT & 0.3 & 0.4 & 99.9 & 93.7 & n/a & n/a \\
\bottomrule
\end{tabular}
\caption{Learned prior against $\mathcal{N}(0,I)$, 50k steps, test split.}
\label{tab:ablation}
\end{table}

The effect of the learned prior depends on the task.
\begin{itemize}
  \item \textbf{Sudoku: more diversity at no cost.} Equally valid (99.5\% vs 99.7\%), equally firm on forced answers
  (1.10 vs 1.11), and more diverse: ratio 0.40 against 0.27, 23.2 against 9.1 distinct valid solutions at 256+.
  \item \textbf{N-Queens and graph colouring: errors, not coverage.} Half as valid or less (48\% vs 97\%, 37\% vs
  99.5\%), with forced answers changed far more often (37 and 7.5 distinct answers per unique input) and the same
  ratio to uniform.
  \item \textbf{Mazes: no real difference} (74.2\% vs 76.7\% valid, ratio 0.15 vs 0.13).
  \item \textbf{Graph colouring's instability destroyed the model before the NaN}: 7.1\% validity at step 41{,}660.
  \item \textbf{Q3 is unchanged}: the unseen share matches a uniform sampler's with the learned prior too (Sudoku
  77.7\% vs 76.8\%, graph colouring 90.0\% vs 89.7\%).
\end{itemize}
We had predicted that $\mathcal{N}(0,I)$ would be as diverse or less, and clearly more valid where the learned prior
was unstable; both held. Our earlier Sudoku prediction, that 30--60\% of valid samples on training inputs would be
unseen, was wrong (77.7\%).

\section{Next steps}
\label{sec:next}

\begin{itemize}
  \item \textbf{150k-step baseline, two seeds, every checkpoint scored on a validation split} (new inputs from the same
  generators, checked for no overlap). The reported number stays the final checkpoint; the curves show whether
  validity and diversity have stopped changing, whether training and validation inputs agree, and whether diversity
  drifts down over long training under an objective indifferent to it. Graph colouring is also trained from scratch
  alongside a continuation of its 50k run, to check that continuing a run is equivalent to one long run.
  \item \textbf{Chance baseline}: the validity of a random answer, $k(x)$ over the size of the answer space, per task
  and group; needed to read SAT and sparse graph colouring.
  \item \textbf{GRAM} as the second baseline, then interventions.
\end{itemize}

\section{Known caveats}

\begin{itemize}
  \item \textbf{One seed} per task. Seed variance on N-Queens was up to 8 points in report v1.
  \item \textbf{SAT's halting target} counts only stored solutions, so its halting head is trained on a stricter
  criterion than validity.
  \item \textbf{Stored solutions bound the error anatomy}: an invalid sample may be close to a solution that is not
  stored.
  \item \textbf{The recombination share needs a uniform reference} per task before it can distinguish generalisation
  from recombination.
  \item \textbf{Fixed step budget}: all runs here train $\sim$50k steps regardless of dataset size; N-Queens has fewer
  training inputs and so more epochs.
\end{itemize}

\appendix
\section{Per-group results}
\label{app:groups}

Generated from \texttt{results/<task>-<split>/summary.md}. Cells off: median cells between an invalid sample and the
nearest stored solution. Blends: share of invalid samples on multi-solution inputs made entirely of cells from stored
solutions. Held-out figures are per group: unseen share of valid samples vs a uniform sampler's.

\begin{table}[H]
\centering\small
\begin{tabular}{@{}lrcccccccc@{}}
\toprule
& & & Any-valid & \multicolumn{3}{c}{Distinct valid@100} & & \multicolumn{2}{c}{Invalid} \\
\cmidrule(lr){5-7}\cmidrule(l){9-10}
Solutions & Inputs & Validity & @100 & model & uniform & ratio & Coverage & cells off & blends \\
\midrule
all & 2000 & 74.2 & 93.5 & 1.8 & 33.6 & 0.15 & 11.0 & 10 & 77.5 \\
1 & 667 & 76.4 & 88.2 & -- & -- & -- & -- & 18 & -- \\
2--3 & 149 & 78.3 & 94.0 & 1.1 & 2.3 & 0.52 & 48.5 & 10 & 24.3 \\
4--15 & 296 & 77.4 & 94.9 & 1.3 & 7.6 & 0.20 & 18.2 & 4 & 63.6 \\
16--63 & 296 & 71.1 & 95.3 & 1.5 & 25.6 & 0.08 & 4.8 & 6 & 89.3 \\
64--255 & 296 & 69.2 & 96.6 & 1.8 & 50.9 & 0.06 & 1.4 & 8 & 88.5 \\
256+ & 296 & 72.1 & 99.0 & 3.0 & 66.3 & 0.06 & 0.6 & 12 & 85.2 \\
\bottomrule
\end{tabular}
\caption{Maze, learned prior, test split. Unique-solution inputs: 11.73 distinct answers per input.}
\end{table}

\begin{table}[H]
\centering\small
\begin{tabular}{@{}lrcccccccc@{}}
\toprule
& & & Any-valid & \multicolumn{3}{c}{Distinct valid@100} & & \multicolumn{2}{c}{Invalid} \\
\cmidrule(lr){5-7}\cmidrule(l){9-10}
Solutions & Inputs & Validity & @100 & model & uniform & ratio & Coverage & cells off & blends \\
\midrule
all & 2000 & 76.7 & 93.2 & 1.7 & 31.9 & 0.13 & 9.9 & 13 & 60.7 \\
1 & 667 & 75.0 & 86.4 & -- & -- & -- & -- & 16 & -- \\
2--3 & 149 & 79.6 & 90.6 & 1.0 & 2.2 & 0.48 & 43.2 & 14 & 14.0 \\
4--15 & 296 & 87.6 & 96.3 & 1.2 & 7.9 & 0.17 & 16.2 & 11 & 36.8 \\
16--63 & 296 & 89.2 & 97.6 & 1.4 & 29.6 & 0.06 & 4.4 & 7 & 68.7 \\
64--255 & 296 & 79.2 & 99.3 & 1.8 & 56.0 & 0.05 & 1.5 & 9 & 81.2 \\
256+ & 296 & 53.2 & 96.6 & 2.7 & 49.0 & 0.10 & 0.6 & 14 & 66.4 \\
\bottomrule
\end{tabular}
\caption{Maze, $\mathcal{N}(0,I)$, test split. Unique-solution inputs: 13.56 distinct answers per input.}
\end{table}

\begin{table}[H]
\centering\small
\begin{tabular}{@{}lrcccccccc@{}}
\toprule
& & & Any-valid & \multicolumn{3}{c}{Distinct valid@100} & & \multicolumn{2}{c}{Invalid} \\
\cmidrule(lr){5-7}\cmidrule(l){9-10}
Solutions & Inputs & Validity & @100 & model & uniform & ratio & Coverage & cells off & blends \\
\midrule
all & 2000 & 74.7 & 94.0 & 1.8 & 33.5 & 0.14 & 10.9 & 11 & 76.8 \\
1 & 660 & 77.3 & 89.1 & -- & -- & -- & -- & 18 & -- \\
2--3 & 144 & 80.3 & 93.1 & 1.1 & 2.2 & 0.49 & 45.6 & 12 & 14.0 \\
4--15 & 300 & 78.2 & 95.7 & 1.4 & 7.6 & 0.22 & 19.8 & 4 & 62.0 \\
16--63 & 299 & 71.1 & 95.3 & 1.4 & 25.3 & 0.08 & 4.9 & 5 & 88.4 \\
64--255 & 297 & 71.3 & 98.0 & 1.9 & 52.3 & 0.05 & 1.6 & 8 & 89.9 \\
256+ & 300 & 69.6 & 98.0 & 2.9 & 63.8 & 0.06 & 0.6 & 13 & 83.6 \\
\bottomrule
\end{tabular}
\caption{Maze, learned prior, first 2000 training inputs. Unique-solution inputs: 12.64 distinct answers per input. Held-out (881 inputs with more than 16 solutions): 77.8\% of valid samples unseen, uniform 76.4\% (16--63: 48.7 vs 45.5, 64--255: 88.5 vs 86.4, 256+: 95.3 vs 96.4); 1.74 distinct unseen per input; 71.3\% of unseen valid samples recombine stored ones.}
\end{table}

\begin{table}[H]
\centering\small
\begin{tabular}{@{}lrcccccccc@{}}
\toprule
& & & Any-valid & \multicolumn{3}{c}{Distinct valid@100} & & \multicolumn{2}{c}{Invalid} \\
\cmidrule(lr){5-7}\cmidrule(l){9-10}
Solutions & Inputs & Validity & @100 & model & uniform & ratio & Coverage & cells off & blends \\
\midrule
all & 2000 & 75.6 & 92.2 & 1.7 & 30.8 & 0.13 & 9.6 & 14 & 58.9 \\
1 & 660 & 74.3 & 82.7 & -- & -- & -- & -- & 18 & -- \\
2--3 & 144 & 77.2 & 91.0 & 1.0 & 2.2 & 0.46 & 41.4 & 14 & 8.1 \\
4--15 & 300 & 86.7 & 95.3 & 1.1 & 7.9 & 0.17 & 16.3 & 11 & 34.5 \\
16--63 & 299 & 90.5 & 98.3 & 1.4 & 28.8 & 0.06 & 4.8 & 7 & 73.9 \\
64--255 & 297 & 76.8 & 97.0 & 1.8 & 54.0 & 0.05 & 1.5 & 10 & 76.5 \\
256+ & 300 & 50.3 & 99.7 & 2.7 & 46.4 & 0.10 & 0.6 & 14 & 65.6 \\
\bottomrule
\end{tabular}
\caption{Maze, $\mathcal{N}(0,I)$, first 2000 training inputs. Unique-solution inputs: 14.04 distinct answers per input. Held-out (881 inputs with more than 16 solutions): 72.0\% of valid samples unseen, uniform 71.9\% (16--63: 44.0 vs 45.3, 64--255: 88.6 vs 85.8, 256+: 94.6 vs 96.2); 1.62 distinct unseen per input; 76.2\% of unseen valid samples recombine stored ones.}
\end{table}

\begin{table}[H]
\centering\small
\begin{tabular}{@{}lrcccccccc@{}}
\toprule
& & & Any-valid & \multicolumn{3}{c}{Distinct valid@100} & & \multicolumn{2}{c}{Invalid} \\
\cmidrule(lr){5-7}\cmidrule(l){9-10}
Solutions & Inputs & Validity & @100 & model & uniform & ratio & Coverage & cells off & blends \\
\midrule
all & 6513 & 48.0 & 98.5 & 3.3 & 3.8 & 0.91 & 90.1 & 7 & 21.0 \\
1 & 2398 & 34.5 & 96.5 & -- & -- & -- & -- & 8 & -- \\
2--3 & 2705 & 53.0 & 99.6 & 2.2 & 2.4 & 0.94 & 93.5 & 7 & 13.8 \\
4--15 & 1358 & 61.4 & 100.0 & 4.9 & 5.9 & 0.85 & 84.5 & 7 & 36.8 \\
16--63 & 43 & 65.0 & 100.0 & 12.0 & 17.9 & 0.67 & 64.8 & 7 & 68.7 \\
64--255 & 9 & 60.4 & 100.0 & 26.8 & 41.8 & 0.64 & 35.2 & 8 & 98.6 \\
\bottomrule
\end{tabular}
\caption{N-Queens, learned prior, test split. Unique-solution inputs: 37.17 distinct answers per input.}
\end{table}

\begin{table}[H]
\centering\small
\begin{tabular}{@{}lrcccccccc@{}}
\toprule
& & & Any-valid & \multicolumn{3}{c}{Distinct valid@100} & & \multicolumn{2}{c}{Invalid} \\
\cmidrule(lr){5-7}\cmidrule(l){9-10}
Solutions & Inputs & Validity & @100 & model & uniform & ratio & Coverage & cells off & blends \\
\midrule
all & 6513 & 96.9 & 100.0 & 3.3 & 3.9 & 0.89 & 89.2 & 7 & 22.7 \\
1 & 2398 & 93.2 & 100.0 & -- & -- & -- & -- & 7 & -- \\
2--3 & 2705 & 98.9 & 100.0 & 2.2 & 2.4 & 0.92 & 92.3 & 7 & 11.8 \\
4--15 & 1358 & 99.4 & 100.0 & 5.0 & 6.0 & 0.84 & 83.8 & 8 & 61.0 \\
16--63 & 43 & 99.8 & 100.0 & 12.9 & 18.5 & 0.71 & 70.6 & 8 & 100.0 \\
64--255 & 9 & 99.9 & 100.0 & 34.3 & 56.1 & 0.61 & 44.1 & 9 & 100.0 \\
\bottomrule
\end{tabular}
\caption{N-Queens, $\mathcal{N}(0,I)$, test split. Unique-solution inputs: 5.17 distinct answers per input.}
\end{table}

\begin{table}[H]
\centering\small
\begin{tabular}{@{}lrcccccccc@{}}
\toprule
& & & Any-valid & \multicolumn{3}{c}{Distinct valid@100} & & \multicolumn{2}{c}{Invalid} \\
\cmidrule(lr){5-7}\cmidrule(l){9-10}
Solutions & Inputs & Validity & @100 & model & uniform & ratio & Coverage & cells off & blends \\
\midrule
all & 2000 & 47.0 & 98.3 & 3.2 & 3.8 & 0.91 & 90.6 & 7 & 20.1 \\
1 & 747 & 33.1 & 96.0 & -- & -- & -- & -- & 8 & -- \\
2--3 & 843 & 53.4 & 99.5 & 2.2 & 2.4 & 0.95 & 94.5 & 7 & 13.1 \\
4--15 & 395 & 59.2 & 100.0 & 4.9 & 6.0 & 0.85 & 83.9 & 7 & 35.2 \\
16--63 & 12 & 59.8 & 100.0 & 12.9 & 19.5 & 0.67 & 59.8 & 7 & 74.5 \\
64--255 & 3 & 59.0 & 100.0 & 23.7 & 42.1 & 0.56 & 28.2 & 9 & 54.5 \\
\bottomrule
\end{tabular}
\caption{N-Queens, learned prior, first 2000 training inputs. Unique-solution inputs: 38.06 distinct answers per input. Held-out (14 inputs with more than 16 solutions): 31.0\% of valid samples unseen, uniform 33.0\% (16--63: 15.7 vs 20.4, 64--255: 88.1 vs 80.2); 6.86 distinct unseen per input; 39.6\% of unseen valid samples recombine stored ones.}
\end{table}

\begin{table}[H]
\centering\small
\begin{tabular}{@{}lrcccccccc@{}}
\toprule
& & & Any-valid & \multicolumn{3}{c}{Distinct valid@100} & & \multicolumn{2}{c}{Invalid} \\
\cmidrule(lr){5-7}\cmidrule(l){9-10}
Solutions & Inputs & Validity & @100 & model & uniform & ratio & Coverage & cells off & blends \\
\midrule
all & 2000 & 96.6 & 100.0 & 3.3 & 3.9 & 0.89 & 89.4 & 7 & 16.2 \\
1 & 747 & 92.3 & 99.9 & -- & -- & -- & -- & 7 & -- \\
2--3 & 843 & 99.0 & 100.0 & 2.2 & 2.4 & 0.93 & 93.3 & 7 & 6.3 \\
4--15 & 395 & 99.3 & 100.0 & 5.0 & 6.0 & 0.82 & 81.9 & 8 & 44.6 \\
16--63 & 12 & 99.8 & 100.0 & 15.8 & 22.7 & 0.73 & 70.9 & 8 & 33.3 \\
64--255 & 3 & 100.0 & 100.0 & 32.3 & 57.7 & 0.57 & 41.1 & -- & -- \\
\bottomrule
\end{tabular}
\caption{N-Queens, $\mathcal{N}(0,I)$, first 2000 training inputs. Unique-solution inputs: 5.53 distinct answers per input. Held-out (14 inputs with more than 16 solutions): 30.6\% of valid samples unseen, uniform 35.7\% (16--63: 16.6 vs 23.5, 64--255: 81.7 vs 80.1); 9.57 distinct unseen per input; 34.7\% of unseen valid samples recombine stored ones.}
\end{table}

\begin{table}[H]
\centering\small
\begin{tabular}{@{}lrcccccccc@{}}
\toprule
& & & Any-valid & \multicolumn{3}{c}{Distinct valid@100} & & \multicolumn{2}{c}{Invalid} \\
\cmidrule(lr){5-7}\cmidrule(l){9-10}
Solutions & Inputs & Validity & @100 & model & uniform & ratio & Coverage & cells off & blends \\
\midrule
all & 2000 & 99.5 & 100.0 & 10.1 & 39.0 & 0.40 & 33.4 & 16 & 31.6 \\
1 & 667 & 99.9 & 100.0 & -- & -- & -- & -- & 21 & -- \\
2--3 & 267 & 99.6 & 99.6 & 1.6 & 2.2 & 0.72 & 71.8 & 24 & 1.0 \\
4--15 & 267 & 99.5 & 100.0 & 3.3 & 7.0 & 0.51 & 50.9 & 15 & 15.4 \\
16--63 & 267 & 99.0 & 100.0 & 7.9 & 29.7 & 0.28 & 26.8 & 15 & 33.8 \\
64--255 & 266 & 99.0 & 100.0 & 14.7 & 67.2 & 0.22 & 12.2 & 17 & 36.9 \\
256+ & 266 & 99.4 & 100.0 & 23.2 & 89.2 & 0.26 & 5.0 & 11 & 51.2 \\
\bottomrule
\end{tabular}
\caption{Sudoku, learned prior, test split. Unique-solution inputs: 1.10 distinct answers per input.}
\end{table}

\begin{table}[H]
\centering\small
\begin{tabular}{@{}lrcccccccc@{}}
\toprule
& & & Any-valid & \multicolumn{3}{c}{Distinct valid@100} & & \multicolumn{2}{c}{Invalid} \\
\cmidrule(lr){5-7}\cmidrule(l){9-10}
Solutions & Inputs & Validity & @100 & model & uniform & ratio & Coverage & cells off & blends \\
\midrule
all & 2000 & 99.7 & 100.0 & 4.6 & 39.1 & 0.27 & 24.0 & 22 & 3.3 \\
1 & 667 & 99.8 & 100.0 & -- & -- & -- & -- & 28 & -- \\
2--3 & 267 & 99.4 & 100.0 & 1.3 & 2.3 & 0.61 & 61.4 & 20 & 1.3 \\
4--15 & 267 & 99.9 & 100.0 & 2.3 & 7.0 & 0.37 & 37.0 & 22 & 20.0 \\
16--63 & 267 & 99.2 & 100.0 & 4.2 & 29.6 & 0.15 & 14.2 & 18 & 1.9 \\
64--255 & 266 & 99.9 & 100.0 & 6.2 & 67.6 & 0.09 & 5.3 & 23 & 7.9 \\
256+ & 266 & 99.8 & 100.0 & 9.1 & 89.6 & 0.10 & 2.0 & 30 & 5.3 \\
\bottomrule
\end{tabular}
\caption{Sudoku, $\mathcal{N}(0,I)$, test split. Unique-solution inputs: 1.11 distinct answers per input.}
\end{table}

\begin{table}[H]
\centering\small
\begin{tabular}{@{}lrcccccccc@{}}
\toprule
& & & Any-valid & \multicolumn{3}{c}{Distinct valid@100} & & \multicolumn{2}{c}{Invalid} \\
\cmidrule(lr){5-7}\cmidrule(l){9-10}
Solutions & Inputs & Validity & @100 & model & uniform & ratio & Coverage & cells off & blends \\
\midrule
all & 2000 & 99.6 & 100.0 & 10.6 & 39.7 & 0.40 & 33.0 & 13 & 40.0 \\
1 & 644 & 100.0 & 100.0 & -- & -- & -- & -- & 24 & -- \\
2--3 & 279 & 99.8 & 100.0 & 1.6 & 2.3 & 0.73 & 72.5 & 14 & 5.0 \\
4--15 & 260 & 99.3 & 100.0 & 3.3 & 7.5 & 0.48 & 47.9 & 14 & 8.7 \\
16--63 & 271 & 99.3 & 100.0 & 8.4 & 30.0 & 0.29 & 27.8 & 12 & 40.4 \\
64--255 & 264 & 99.3 & 100.0 & 14.8 & 67.6 & 0.22 & 11.9 & 8 & 55.9 \\
256+ & 282 & 99.1 & 100.0 & 24.4 & 89.5 & 0.27 & 4.9 & 12 & 58.1 \\
\bottomrule
\end{tabular}
\caption{Sudoku, learned prior, first 2000 training inputs. Unique-solution inputs: 1.01 distinct answers per input. Held-out (801 inputs with more than 16 solutions): 77.7\% of valid samples unseen, uniform 76.8\% (16--63: 48.9 vs 45.4, 64--255: 85.6 vs 86.0, 256+: 96.4 vs 96.6); 13.88 distinct unseen per input; 58.7\% of unseen valid samples recombine stored ones.}
\end{table}

\begin{table}[H]
\centering\small
\begin{tabular}{@{}lrcccccccc@{}}
\toprule
& & & Any-valid & \multicolumn{3}{c}{Distinct valid@100} & & \multicolumn{2}{c}{Invalid} \\
\cmidrule(lr){5-7}\cmidrule(l){9-10}
Solutions & Inputs & Validity & @100 & model & uniform & ratio & Coverage & cells off & blends \\
\midrule
all & 2000 & 99.9 & 100.0 & 4.8 & 39.9 & 0.26 & 23.6 & 19 & 4.3 \\
1 & 644 & 100.0 & 100.0 & -- & -- & -- & -- & 20 & -- \\
2--3 & 279 & 99.7 & 100.0 & 1.4 & 2.3 & 0.64 & 63.5 & 19 & 0.0 \\
4--15 & 260 & 99.7 & 100.0 & 2.2 & 7.5 & 0.33 & 33.0 & 18 & 2.4 \\
16--63 & 271 & 99.8 & 100.0 & 4.2 & 30.0 & 0.15 & 14.1 & 18 & 6.2 \\
64--255 & 264 & 99.9 & 100.0 & 6.4 & 67.9 & 0.09 & 5.2 & 24 & 8.0 \\
256+ & 282 & 100.0 & 100.0 & 9.5 & 90.2 & 0.11 & 2.0 & 17 & 42.9 \\
\bottomrule
\end{tabular}
\caption{Sudoku, $\mathcal{N}(0,I)$, first 2000 training inputs. Unique-solution inputs: 1.01 distinct answers per input. Held-out (801 inputs with more than 16 solutions): 76.9\% of valid samples unseen, uniform 76.8\% (16--63: 46.1 vs 45.3, 64--255: 85.0 vs 86.0, 256+: 97.1 vs 96.6); 5.71 distinct unseen per input; 58.2\% of unseen valid samples recombine stored ones.}
\end{table}

\begin{table}[H]
\centering\small
\begin{tabular}{@{}lrcccccccc@{}}
\toprule
& & & Any-valid & \multicolumn{3}{c}{Distinct valid@100} & & \multicolumn{2}{c}{Invalid} \\
\cmidrule(lr){5-7}\cmidrule(l){9-10}
Solutions & Inputs & Validity & @100 & model & uniform & ratio & Coverage & cells off & blends \\
\midrule
all & 2000 & 37.0 & 93.3 & 16.9 & 22.2 & 0.74 & 25.0 & 2 & 68.8 \\
1 & 667 & 45.1 & 91.6 & -- & -- & -- & -- & 1 & -- \\
2--3 & 223 & 27.1 & 83.4 & 1.4 & 1.7 & 0.80 & 63.5 & 2 & 16.1 \\
4--15 & 222 & 23.8 & 89.2 & 3.2 & 5.0 & 0.66 & 45.6 & 2 & 46.2 \\
16--63 & 222 & 26.8 & 94.1 & 6.8 & 13.6 & 0.59 & 21.9 & 2 & 78.6 \\
64--255 & 222 & 28.4 & 98.6 & 13.9 & 23.5 & 0.68 & 11.3 & 2 & 92.8 \\
256+ & 444 & 45.9 & 100.0 & 38.0 & 44.7 & 0.86 & 3.8 & 3 & 97.9 \\
\bottomrule
\end{tabular}
\caption{Graph colouring, learned prior, step 37{,}494, test split. Unique-solution inputs: 7.45 distinct answers per input. Valid up to relabelling: 37.0\%.}
\end{table}

\begin{table}[H]
\centering\small
\begin{tabular}{@{}lrcccccccc@{}}
\toprule
& & & Any-valid & \multicolumn{3}{c}{Distinct valid@100} & & \multicolumn{2}{c}{Invalid} \\
\cmidrule(lr){5-7}\cmidrule(l){9-10}
Solutions & Inputs & Validity & @100 & model & uniform & ratio & Coverage & cells off & blends \\
\midrule
all & 2000 & 7.1 & 77.5 & 6.3 & 7.5 & 0.84 & 10.4 & 4 & 46.3 \\
1 & 667 & 5.3 & 73.9 & -- & -- & -- & -- & 5 & -- \\
2--3 & 223 & 2.9 & 54.7 & 0.7 & 1.0 & 0.72 & 33.2 & 4 & 3.1 \\
4--15 & 222 & 2.2 & 63.1 & 1.1 & 1.5 & 0.71 & 17.4 & 4 & 14.2 \\
16--63 & 222 & 2.7 & 71.6 & 1.9 & 2.5 & 0.79 & 6.2 & 3 & 41.8 \\
64--255 & 222 & 4.6 & 86.9 & 3.6 & 4.5 & 0.88 & 3.0 & 3 & 63.3 \\
256+ & 444 & 18.0 & 99.5 & 15.3 & 17.8 & 0.90 & 1.2 & 3 & 84.0 \\
\bottomrule
\end{tabular}
\caption{Graph colouring, learned prior, step 41{,}660, test split. Unique-solution inputs: 69.15 distinct answers per input. Valid up to relabelling: 7.5\%.}
\end{table}

\begin{table}[H]
\centering\small
\begin{tabular}{@{}lrcccccccc@{}}
\toprule
& & & Any-valid & \multicolumn{3}{c}{Distinct valid@100} & & \multicolumn{2}{c}{Invalid} \\
\cmidrule(lr){5-7}\cmidrule(l){9-10}
Solutions & Inputs & Validity & @100 & model & uniform & ratio & Coverage & cells off & blends \\
\midrule
all & 2000 & 99.5 & 100.0 & 35.7 & 48.9 & 0.75 & 46.8 & 1 & 94.5 \\
1 & 667 & 100.0 & 100.0 & -- & -- & -- & -- & 1 & -- \\
2--3 & 223 & 99.7 & 100.0 & 2.1 & 2.3 & 0.94 & 93.7 & 1 & 72.1 \\
4--15 & 222 & 99.2 & 100.0 & 5.8 & 7.4 & 0.81 & 81.0 & 1 & 94.7 \\
16--63 & 222 & 98.8 & 100.0 & 17.4 & 30.2 & 0.59 & 55.8 & 1 & 98.5 \\
64--255 & 222 & 98.7 & 100.0 & 39.3 & 67.3 & 0.58 & 32.1 & 2 & 94.9 \\
256+ & 444 & 99.6 & 100.0 & 74.8 & 93.4 & 0.80 & 8.8 & 3 & 95.1 \\
\bottomrule
\end{tabular}
\caption{Graph colouring, $\mathcal{N}(0,I)$, test split. Unique-solution inputs: 1.02 distinct answers per input. Valid up to relabelling: 99.5\%.}
\end{table}

\begin{table}[H]
\centering\small
\begin{tabular}{@{}lrcccccccc@{}}
\toprule
& & & Any-valid & \multicolumn{3}{c}{Distinct valid@100} & & \multicolumn{2}{c}{Invalid} \\
\cmidrule(lr){5-7}\cmidrule(l){9-10}
Solutions & Inputs & Validity & @100 & model & uniform & ratio & Coverage & cells off & blends \\
\midrule
all & 2000 & 36.4 & 93.9 & 17.1 & 22.5 & 0.74 & 24.0 & 2 & 68.2 \\
1 & 649 & 44.6 & 92.1 & -- & -- & -- & -- & 1 & -- \\
2--3 & 226 & 28.4 & 88.5 & 1.4 & 1.8 & 0.80 & 66.2 & 2 & 15.8 \\
4--15 & 225 & 21.3 & 84.4 & 2.7 & 4.7 & 0.60 & 38.4 & 2 & 40.2 \\
16--63 & 210 & 20.7 & 95.2 & 6.5 & 12.2 & 0.62 & 20.5 & 2 & 79.5 \\
64--255 & 213 & 32.4 & 100.0 & 15.2 & 26.0 & 0.66 & 12.4 & 2 & 93.5 \\
256+ & 477 & 44.9 & 100.0 & 36.9 & 43.6 & 0.86 & 3.9 & 3 & 98.2 \\
\bottomrule
\end{tabular}
\caption{Graph colouring, learned prior, step 37{,}494, first 2000 training inputs. Unique-solution inputs: 7.47 distinct answers per input. Valid up to relabelling: 36.4\%. Held-out (876 inputs with more than 16 solutions): 90.0\% of valid samples unseen, uniform 89.7\% (16--63: 54.1 vs 52.4, 64--255: 85.5 vs 85.7, 256+: 98.3 vs 98.1); 23.71 distinct unseen per input; 98.3\% of unseen valid samples recombine stored ones.}
\end{table}

\begin{table}[H]
\centering\small
\begin{tabular}{@{}lrcccccccc@{}}
\toprule
& & & Any-valid & \multicolumn{3}{c}{Distinct valid@100} & & \multicolumn{2}{c}{Invalid} \\
\cmidrule(lr){5-7}\cmidrule(l){9-10}
Solutions & Inputs & Validity & @100 & model & uniform & ratio & Coverage & cells off & blends \\
\midrule
all & 2000 & 99.4 & 100.0 & 36.5 & 49.9 & 0.76 & 46.2 & 1 & 94.5 \\
1 & 649 & 99.9 & 100.0 & -- & -- & -- & -- & 1 & -- \\
2--3 & 226 & 99.7 & 100.0 & 2.0 & 2.2 & 0.93 & 92.9 & 1 & 65.3 \\
4--15 & 225 & 99.4 & 100.0 & 5.9 & 7.5 & 0.81 & 81.1 & 1 & 98.6 \\
16--63 & 210 & 98.2 & 100.0 & 18.3 & 30.7 & 0.61 & 57.4 & 1 & 99.0 \\
64--255 & 213 & 98.6 & 100.0 & 39.4 & 67.2 & 0.58 & 32.2 & 2 & 92.2 \\
256+ & 477 & 99.5 & 100.0 & 73.9 & 93.2 & 0.79 & 9.0 & 3 & 97.0 \\
\bottomrule
\end{tabular}
\caption{Graph colouring, $\mathcal{N}(0,I)$, first 2000 training inputs. Unique-solution inputs: 1.03 distinct answers per input. Valid up to relabelling: 99.4\%. Held-out (876 inputs with more than 16 solutions): 84.7\% of valid samples unseen, uniform 84.7\% (16--63: 50.0 vs 49.2, 64--255: 85.4 vs 86.0, 256+: 97.8 vs 97.8); 49.97 distinct unseen per input; 98.5\% of unseen valid samples recombine stored ones.}
\end{table}

\begin{table}[H]
\centering\small
\begin{tabular}{@{}lrcccccccc@{}}
\toprule
& & & Any-valid & \multicolumn{3}{c}{Distinct valid@100} & & \multicolumn{2}{c}{Invalid} \\
\cmidrule(lr){5-7}\cmidrule(l){9-10}
Solutions & Inputs & Validity & @100 & model & uniform & ratio & Coverage & cells off & blends \\
\midrule
all & 2000 & 0.3 & 12.0 & 0.4 & 0.4 & 1.00 & 0.1 & 6 & 13.5 \\
1 & 667 & 0.0 & 0.3 & -- & -- & -- & -- & 8 & -- \\
2--3 & 267 & 0.0 & 0.0 & 0.0 & 0.0 & -- & 0.0 & 7 & 0.0 \\
4--15 & 267 & 0.0 & 3.0 & 0.0 & 0.0 & -- & 0.3 & 6 & 2.5 \\
16--63 & 267 & 0.0 & 4.1 & 0.0 & 0.0 & -- & 0.1 & 6 & 8.3 \\
64--255 & 266 & 0.2 & 15.4 & 0.2 & 0.2 & 1.00 & 0.1 & 5 & 19.5 \\
256+ & 266 & 1.9 & 67.3 & 1.9 & 1.9 & 1.00 & 0.1 & 5 & 37.8 \\
\bottomrule
\end{tabular}
\caption{3-SAT, learned prior, test split. Unique-solution inputs: 99.93 distinct answers per input.}
\end{table}

\begin{table}[H]
\centering\small
\begin{tabular}{@{}lrcccccccc@{}}
\toprule
& & & Any-valid & \multicolumn{3}{c}{Distinct valid@100} & & \multicolumn{2}{c}{Invalid} \\
\cmidrule(lr){5-7}\cmidrule(l){9-10}
Solutions & Inputs & Validity & @100 & model & uniform & ratio & Coverage & cells off & blends \\
\midrule
all & 2000 & 0.4 & 12.9 & 0.6 & 0.7 & 0.90 & 0.5 & 6 & 22.0 \\
1 & 667 & 0.0 & 0.1 & -- & -- & -- & -- & 7 & -- \\
2--3 & 267 & 0.0 & 1.1 & 0.0 & 0.0 & 1.33 & 0.7 & 6 & 0.1 \\
4--15 & 267 & 0.0 & 3.4 & 0.0 & 0.0 & 0.54 & 0.5 & 6 & 3.4 \\
16--63 & 267 & 0.2 & 12.4 & 0.2 & 0.2 & 0.91 & 0.6 & 5 & 14.6 \\
64--255 & 266 & 0.7 & 25.9 & 0.6 & 0.7 & 0.91 & 0.4 & 5 & 33.5 \\
256+ & 266 & 2.5 & 53.8 & 2.0 & 2.5 & 0.89 & 0.2 & 5 & 59.7 \\
\bottomrule
\end{tabular}
\caption{3-SAT, $\mathcal{N}(0,I)$, test split. Unique-solution inputs: 93.68 distinct answers per input.}
\end{table}

\begin{table}[H]
\centering\small
\begin{tabular}{@{}lrcccccccc@{}}
\toprule
& & & Any-valid & \multicolumn{3}{c}{Distinct valid@100} & & \multicolumn{2}{c}{Invalid} \\
\cmidrule(lr){5-7}\cmidrule(l){9-10}
Solutions & Inputs & Validity & @100 & model & uniform & ratio & Coverage & cells off & blends \\
\midrule
all & 2000 & 0.3 & 13.4 & 0.5 & 0.5 & 1.00 & 0.2 & 6 & 13.2 \\
1 & 657 & 0.0 & 0.0 & -- & -- & -- & -- & 8 & -- \\
2--3 & 269 & 0.0 & 0.4 & 0.0 & 0.0 & -- & 0.2 & 7 & 0.0 \\
4--15 & 277 & 0.0 & 3.6 & 0.0 & 0.0 & -- & 0.5 & 6 & 1.6 \\
16--63 & 274 & 0.1 & 5.1 & 0.1 & 0.1 & 1.02 & 0.2 & 6 & 8.9 \\
64--255 & 248 & 0.2 & 18.1 & 0.2 & 0.2 & 1.00 & 0.1 & 5 & 15.6 \\
256+ & 275 & 2.0 & 72.0 & 2.0 & 2.0 & 1.00 & 0.1 & 5 & 40.2 \\
\bottomrule
\end{tabular}
\caption{3-SAT, learned prior, first 2000 training inputs. Unique-solution inputs: 99.93 distinct answers per input. Held-out (781 inputs with more than 16 solutions): 97.5\% of valid samples unseen, uniform 96.6\% (16--63: 46.7 vs 47.0, 64--255: 94.0 vs 88.4, 256+: 99.3 vs 98.8); 0.75 distinct unseen per input; 98.3\% of unseen valid samples recombine stored ones.}
\end{table}

\begin{table}[H]
\centering\small
\begin{tabular}{@{}lrcccccccc@{}}
\toprule
& & & Any-valid & \multicolumn{3}{c}{Distinct valid@100} & & \multicolumn{2}{c}{Invalid} \\
\cmidrule(lr){5-7}\cmidrule(l){9-10}
Solutions & Inputs & Validity & @100 & model & uniform & ratio & Coverage & cells off & blends \\
\midrule
all & 2000 & 0.7 & 14.6 & 0.8 & 1.0 & 0.89 & 0.5 & 6 & 21.8 \\
1 & 657 & 0.0 & 0.3 & -- & -- & -- & -- & 7 & -- \\
2--3 & 269 & 0.0 & 1.5 & 0.0 & 0.0 & -- & 0.6 & 7 & 0.1 \\
4--15 & 277 & 0.1 & 4.3 & 0.1 & 0.1 & 0.90 & 0.7 & 6 & 2.7 \\
16--63 & 274 & 0.2 & 7.7 & 0.2 & 0.2 & 0.86 & 0.4 & 5 & 14.2 \\
64--255 & 248 & 0.8 & 31.5 & 0.7 & 0.8 & 0.92 & 0.5 & 5 & 32.5 \\
256+ & 275 & 3.7 & 63.6 & 2.9 & 3.7 & 0.88 & 0.3 & 5 & 61.9 \\
\bottomrule
\end{tabular}
\caption{3-SAT, $\mathcal{N}(0,I)$, first 2000 training inputs. Unique-solution inputs: 93.85 distinct answers per input. Held-out (781 inputs with more than 16 solutions): 96.5\% of valid samples unseen, uniform 95.4\% (16--63: 76.9 vs 64.3, 64--255: 94.7 vs 88.9, 256+: 97.9 vs 98.3); 1.27 distinct unseen per input; 97.3\% of unseen valid samples recombine stored ones.}
\end{table}


\clearpage

\begin{thebibliography}{99}\small
\bibitem{baek2026gram} J.~Baek, M.~Jo, M.~Kim, M.~Ren, Y.~Bengio, and S.~Ahn. Generative recursive reasoning.
\emph{arXiv:2605.19376}, 2026.
\bibitem{jolicoeur2025trm} A.~Jolicoeur-Martineau. Less is more: Recursive reasoning with tiny networks.
\emph{arXiv:2510.04871}, 2025.
\bibitem{nandwani2021oneofmany} Y.~Nandwani, D.~Jindal, Mausam, and P.~Singla. Neural learning of one-of-many
solutions for combinatorial problems in structured output spaces. In \emph{ICLR}, 2021.
\bibitem{bhattacharyya2018bestofmany} A.~Bhattacharyya, B.~Schiele, and M.~Fritz. Accurate and diverse sampling of
sequences based on a ``best of many'' sample objective. In \emph{CVPR}, 2018.
\bibitem{li2018imle} K.~Li and J.~Malik. Implicit maximum likelihood estimation. \emph{arXiv:1809.09087}, 2018.
\end{thebibliography}

\end{document}
